%Generazione delle variabili che andranno a sostituire quelle del template `HomePage.tex`
\newcommand{\documento}{Manuale del Manutentore}
\newcommand{\nomedocumentofisico}{ManualeManutentore\_v1\_0\_0.pdf}
\newcommand{\redazione}{\SM}
\newcommand{\verifica}{\GN\\ & \MV}
\newcommand{\approvazione}{\MP}
\newcommand{\uso}{Esterno}
\newcommand{\destinateTo}{\TV, \\ & \RC, \\ & \ZU}
\newcommand{\datacreazione}{09 Maggio 2016}
\newcommand{\datamodifica}{09 Maggio 2016}
\newcommand{\stato}{In redazione}
\newcommand{\versione}{1}

\def\TABELLE{false}	%abilita - disabilita l'indice delle tabelle
\def\FIGURE{false} 	%abilita - disabilita l'indice delle figure

%Layout del documento 
\documentclass[a4paper,11pt]{article}

%***IMPORTAZIONE PACKAGE***
\usepackage{ifthen}
\usepackage[italian]{babel}
\usepackage[utf8]{inputenc}
\usepackage[T1]{fontenc}
\usepackage{float}
\usepackage{chapterbib}
\usepackage{graphicx}
\usepackage[a4paper,top=2.5cm,bottom=2.5cm,left=2.5cm,right=2.5cm]{geometry}
\usepackage[colorlinks=true, urlcolor=black, citecolor=black, linkcolor=black]{hyperref}
\usepackage{booktabs}
\usepackage{fancyhdr}
\usepackage{totpages}
\usepackage{tabularx, array}
\usepackage{dcolumn}
\usepackage{epstopdf}
\usepackage{booktabs}
\usepackage{fancyhdr}
\usepackage{longtable}
\usepackage{calc}
\usepackage{datatool}
\usepackage[bottom]{footmisc}
\usepackage{listings} 
\usepackage{textcomp}
\usepackage{titlesec}
\usepackage{rotating} 
\usepackage{multirow}
\usepackage{placeins}
\usepackage{color}
\usepackage[table,usenames,dvipsnames]{xcolor}
\usepackage{hyperref}
\usepackage{makecell}
\usepackage{hyperref}


%***STILE PAGINA***
\pagestyle{fancy}
%no indentazione paragrafo
\setlength{\parindent}{0pt}

%***INTESTAZIONE***
\lhead{\Large{\progetto} \\ \footnotesize{\documento}}
\rhead{\includegraphics[keepaspectratio = true, width = 25px] {../../Template/icone/LogoGruppo.png}}
\renewcommand{\headrulewidth}{0.4pt}  %Linea sotto l'intestazione

%***PIÈ DI PAGINA***
\lfoot{\textit{\gruppoLink}\\ \footnotesize{\email}}
\rfoot{\thepage} %per le prime pagine: mostra solo il numero romano
\cfoot{}
\renewcommand{\footrulewidth}{0.4pt}   %Linea sopra il piè di pagina

%***INSERIMENTO DI NUOVE SOTTOSEZIONI
\setcounter{secnumdepth}{7}		% mostra nel documento fino al livello 8 (1.2.3.4.5.6.7.8)
\setcounter{tocdepth}{7}			% mostra nell'indice fino al livello 8 (1.2.3.4.5.6.7.8)

%***LA SOTTOSEZIONE PARAGRAPH VIENE VISUALIZZATA COME UNA SECTION
\titleformat{\paragraph}{\normalfont\normalsize\bfseries}{\theparagraph}{1em}{}
\titlespacing*{\paragraph}{0pt}{3.25ex plus 1ex minus .2ex}{1.5ex plus .2ex}

\titleformat{\subparagraph}{\normalfont\normalsize\bfseries}{\thesubparagraph}{1em}{}
\titlespacing*{\subparagraph}{0pt}{3.25ex plus 1ex minus .2ex}{1.5ex plus .2ex}

\makeatletter
\newcounter{subsubparagraph}[subparagraph]
\renewcommand\thesubsubparagraph{%
  \thesubparagraph.\@arabic\c@subsubparagraph}
\newcommand\subsubparagraph{%
  \@startsection{subsubparagraph}    % counter
    {6}                              % level
    {\parindent}                     % indent
    {3.25ex \@plus 1ex \@minus .2ex} % beforeskip
    {0.75em}                           % afterskip
    {\normalfont\normalsize\bfseries}}
\newcommand\l@subsubparagraph{\@dottedtocline{6}{10em}{5.5em}} %gestione dell'indice
\newcommand{\subsubparagraphmark}[1]{}
\makeatother

\makeatletter
\newcounter{subsubsubparagraph}[subsubparagraph]
\renewcommand\thesubsubsubparagraph{%
  \thesubsubparagraph.\@arabic\c@subsubsubparagraph}
\newcommand\subsubsubparagraph{%
  \@startsection{subsubsubparagraph}    % counter
    {7}                              % level
    {\parindent}                     % indent
    {3.25ex \@plus 1ex \@minus .2ex} % beforeskip
    {0.75em}                           % afterskip
    {\normalfont\normalsize\bfseries}}
\newcommand\l@subsubsubparagraph{\@dottedtocline{7}{10em}{6.5em}} %gestione dell'indice
\newcommand{\subsubsubparagraphmark}[1]{}
\makeatother


%Comandi generali
%Generali
\newcommand{\progetto}{QuizziPedia}
\newcommand{\gruppo}{TheFellowshipOfTheCode}
\newcommand{\gruppoLink}{\href{http://thefellowshipofthecode.github.io/}{TheFellowshipOfTheCode}}
\newcommand{\email}{\href{mailto:thefellowshipofthecode@gmail.com}{thefellowshipofthecode@gmail.com}}

%Documenti
\newcommand{\AdR}{Analisi dei Requisiti}
\newcommand{\NdP}{Norme di Progetto}
\newcommand{\PdP}{Piano di Progetto}
\newcommand{\SdF}{Studio di Fattibilità}
\newcommand{\PdQ}{Piano di Qualifica}
\newcommand{\VI}{Verbale Interno}
\newcommand{\VE}{Verbale Esterno}
\newcommand{\ST}{Specifica Tecnica}
\newcommand{\DDP}{Definizione di Prodotto}
\newcommand{\MU}{Manuale Utente}
\newcommand{\G}{Glossario}
\newcommand{\LdP}{Lettera di Presentazione}
\newcommand{\NdPv}{NormeDiProgetto\_v\_1\_0\_0}
\newcommand{\PdPv}{PianoDiProgetto\_v\_1\_0\_0}
\newcommand{\PdQv}{PianoDiQualifica\_v\_1\_0\_0}
\newcommand{\SdFv}{StudioDiFattibilità\_v\_1\_0\_0}

%Componenti del gruppo
\newcommand{\AF}{Alberto Ferrara}
\newcommand{\SM}{Simone Magagna}
\newcommand{\FB}{Franco Berton}
\newcommand{\MP}{Marco Prelaz}
\newcommand{\MV}{Mattia Varotto}
\newcommand{\GN}{Matteo Gnoato}
\newcommand{\GR}{Matteo Granzotto}

%Ruoli
\newcommand{\RdP}{Responsabile di Progetto}
\newcommand{\Res}{Responsabile}
\newcommand{\Amm}{Amministratore}
\newcommand{\Ver}{Verificatore}
\newcommand{\Prog}{Progettista}
\newcommand{\Progr}{Programmatore}
\newcommand{\Ana}{Analista}
\newcommand{\RdPs}{Responsabili di Progetto}
\newcommand{\Ress}{Responsabile}
\newcommand{\Amms}{Amministratori}
\newcommand{\Vers}{Verificatori}
\newcommand{\Progs}{Progettisti}
\newcommand{\Progrs}{Programmatori}
\newcommand{\Anas}{Analisti}

%Professori e proponente
\newcommand{\TV}{Prof. Tullio Vardanega}
\newcommand{\RC}{Prof. Riccardo Cardin}
\newcommand{\ZU}{Zucchetti S.P.A.}
\newcommand{\proponente}{Zucchetti S.P.A.}

\newcommand{\diaryEntry}[5]{#2 & \emph{#4} & #3 & #5 & #1\\ \hline}

%comando per una nuova riga nella tabella del diario delle modifiche
\newcommand{\specialcell}[2][c]{%
	\begin{tabular}[#1]{@{}c@{}}#2\end{tabular}}

\renewcommand*\sectionmark[1]{\markboth{#1}{}}
\renewcommand*\subsectionmark[1]{\markright{#1}}

%Pediodi di lavoro 
\newcommand{\AR}{Analisi dei Requisiti}
\newcommand{\AD}{Analisi dei Requisiti in Dettaglio}
\newcommand{\PA}{Progettazione Architetturale}
\newcommand{\PD}{Progettazione di Dettaglio}
\newcommand{\CO}{Codifica}
\newcommand{\VV}{Verifica e Validazione}

% Revisioni
\newcommand{\RR}{Revisione dei Requisiti}
\newcommand{\RP}{Revisione di Progettazione}
\newcommand{\RQ}{Revisione di Qualifica}
\newcommand{\RA}{Revisione di Accettazione}

% Comandi analisi dei requisiti
\newcommand{\uau}{utente autenticato}
\newcommand{\uaus}{utenti autenticati}
\newcommand{\uaupro}{utente autenticato pro}
\newcommand{\uauspro}{utenti autenticati pro}

\newcommand{\myincludegraphics}[2][]{%
	\setbox0=\hbox{\phantom{X}}%
	\vtop{
		\hbox{\phantom{X}}
		\vskip-\ht0
		\hbox{\includegraphics[#1]{#2}}}}

\begin{document}
	
%inclusione template HomePage
\begin{center}

%\includegraphics[width=1em]{../../../Template/icone/LogoGruppo.png}
\begin{large} \textbf{\gruppoLink} \end{large}
%\includegraphics[width=1em]{../../../Template/icone/LogoGruppo.png}
\vspace{0.2em}

\hrule
\vspace{3em}

\includegraphics[keepaspectratio = true, width=8cm]{../../../Template/icone/LogoGruppo.png}

%Prima pagina senza intestazione né piè di pagina	
\thispagestyle{empty}

%Le informazioni del documento sono ancorate a fine pagina
\vfill

%Copertina
\begin{center} 
  \begin{Huge}
  {\fontsize{15mm}{20mm}\selectfont \progetto} 
  \end{Huge}
\end{center}

\begin{Huge} \documento \end{Huge}

\begin{center}
\textbf{Informazioni sul documento} \\ \vspace{2em}
\small
\begin{tabular}{r|l}
	\textbf{Nome Documento} & \nomedocumentofisico \\
	\textbf{Versione}	& 1\\
	\textbf{Data di Creazione} & \datacreazione\\
	\textbf{Data ultima modifica} & \datamodifica\\
	\textbf{Stato} & \stato \\
	\textbf{Redazione}	& \redazione\\
	\textbf{Verifica}	& \verifica\\
	\textbf{Approvazione}	& \approvazione\\
	\textbf{Uso}  & \uso\\
	\textbf{Distribuzione} & \gruppo \\
	\textbf{Destinato a}  &  \destinateTo \\
	\textbf{Email di riferimento} & \email
\end{tabular}
\end{center}

\normalsize
%Sommario
\textbf{Sommario\\} 
Documento contenente le norme di progetto che il gruppo \textit{\gruppo} seguirà durante tutte le fasi di realizzazione del prodotto \textit{\progetto}.

%\vfill %cosa fa?
\end{center}
\clearpage


%Registro delle modifiche e indice 
%si usa la numerazione romana per gli indici e la tabella delle modifiche
\pagenumbering{Roman}
\newcommand{\modificheuno} 
{	
	0.0.3 & Inseriti i primi riferimenti informativi & \specialcell[t]{\GN \\ \prog} & 2016-02-26
	\\\midrule
	0.0.2 & Stesura scopo del documento, scopo del prodotto , riferimenti normativi & \specialcell[t]{\GN \\ \prog} & 2016-02-26
	\\\midrule
	0.0.1 & Creato template documento & \specialcell[t]{\GN \\ \prog} & 2016-02-26
	\\\midrule
}
\newcommand{\modifichedue}
{
}
%Inserisce il link all'indice
%\addcontentsline{toc}{section}{Indice}
\tableofcontents
\clearpage 

%Se è stata impostata a true la variabile per la lista delle tabelle, la mostra
\ifthenelse{\equal{\TABELLE}{true}} 
{\listoftables \newpage}{}

%Se è stata impostata a true la variabile per la lista delle figure, la mostra
\ifthenelse{\equal{\FIGURE}{true}}
{\listoffigures \newpage}{}

%Da qui comincia la numerazione normale
\pagenumbering{arabic}

%Imposta il formato di visualizzazione
\rfoot{\thepage~di~\pageref{TotPages}}
\newpage
\listoffigures
\newpage
\listoftables
\newpage
%sezioni documento
\newpage
\section{Introduzione}

\subsection{Scopo del documento}
Il presente documento ha lo scopo di definire in dettaglio la struttura e il funzionamento delle componenti del progetto \progetto. Questo documento servirà come guida per i \textit{\Progrs} del gruppo \gruppo fornendo direttive e vincoli per la realizzazione del \textit{progetto\ped{G}}.

\subsection{Scopo del prodotto}
Lo scopo del prodotto è di permettere la creazione e gestione di questionari in grado di identificare le lacune dei candidati prima, durante e al termine di un corso di formazione. 
\\Il sistema dovrà offrire le seguenti funzionalità:
\begin{itemize}
	\item
	Archiviare questionari in un server suddivisi per argomento;
	\item
	Somministrare all'utente, tramite un'interfaccia, questionari specifici per argomento scelto;
	\item
	Verificare e valutare i questionari scelti dagli utenti in base alle risposte date.
\end{itemize}
La parte destinata ai creatori di questionari dovrà essere fruibile attraverso un \textit{browser\ped{G}} desktop, abilitato all'utilizzo delle tecnologie \textit{HTML5\ped{G}}, \textit{CSS3\ped{G}} e \textit{JavaScript\ped{G}}. La parte destinata agli esaminandi sarà utilizzabile su qualunque dispositivo: dal personal computer ai tablet e smartphone.

\subsection{Glossario}
Al fine di evitare ogni ambiguità i termini tecnici del dominio del progetto, gli acronimi e le parole che necessitano di ulteriori spiegazioni saranno nei vari documenti marcate con il pedice \ped{G} e quindi presenti nel documento \textit{\G}.


\subsection{Riferimenti}
\subsubsection{Normativi}
\begin{itemize}
	\item \textit{\NdPv};
	\item \textit{\AdRvDue};
\end{itemize}
\subsubsection{Informativi}
\begin{itemize}
	\item \textbf{Ingegneria del software - Ian Sommerville - 8a edizione (2007)}: \\
	Parte terza: Progettazione, capitolo 11: Progettazione architetturale, Capitolo 14: Progettazione orientata agli oggetti;
	\item \textbf{Design Patterns} - Erich Gamma, Richard Helm, Ralph Johnson, John Vlissides - 1a edizione italiana (2006);
	\item \textbf{Slide dell'insegnamento - Design patterns:}
	\begin{itemize}
		\item Strutturali: \url{http://www.math.unipd.it/~tullio/IS-1/2015/Dispense/E07.pdf};
		\item Creazionali: \url{http://www.math.unipd.it/~tullio/IS-1/2015/Dispense/E08.pdf}
		\item Comportamentali: \url{http://www.math.unipd.it/~tullio/IS-1/2015/Dispense/E09.pdf}
		\item Architetturali:
			\begin{itemize}
				\item \url{http://www.math.unipd.it/~rcardin/sweb/Design%20Pattern%20Architetturali%20-%20Model%20View%20Controller_4x4.pdf};
				\item \url{http://www.math.unipd.it/~rcardin/sweb/Design%20Pattern%20Architetturali%20-%20Dependency%20Injection_4x4.pdf}.
			\end{itemize} 
	\end{itemize}
	\item \textbf{Martin Fowler - UML\ped{G} Distilled} - 2nd edition;
	\item \textbf{Slide dell'insegnamento - Diagrammi delle classi}: \\
		\url{http://www.math.unipd.it/~tullio/IS-1/2015/Dispense/E03.pdf}
	\item \textbf{Slide dell'insegnamento - Diagrammi dei packages:} \\
		\url{http://www.math.unipd.it/~tullio/IS-1/2015/Dispense/E04.pdf}
	\item \textbf{Slide dell'insegnamento - Diagrammi di sequenza:} \\
		\url{http://www.math.unipd.it/~tullio/IS-1/2015/Dispense/E05.pdf}
	\item \textbf{Documentazione del \textit{Framework\ped{G}MEAN\ped{G}.js}:} \\
		\url{http://learn.mean.io/}
	\item \textbf{Documentazione della \textit{piattaforma} Node.js:} \\
		\url{https://nodejs.org/api/}
	\item \textbf{Giuda all'utilizzo dei middleware Express:} \\
		\url{http://expressjs.com/it/guide/using-middleware.html}
	\item \textbf{Guida all'utilizzo dei middleware Passport:} \\
		\url{http://passportjs.org/docs}
	\item \textbf{Manuale del database \textit{MongoDB\ped{G}}:} \\
		\url{https://docs.mongodb.org/manual/}
	\item \textbf{Documentazione dell'interfaccia REST:}
		\begin{itemize}
			\item \textit{Descrizione di REST:} \url{https://it.wikipedia.org/wiki/Representational_State_Transfer}
			\item \textit{Descrizione risorse REST:} \url{http://stashboard.readthedocs.org/en/latest/restapi.html}
		\end{itemize}
	\item \textbf{Documentazione del \textit{framework\ped{G} AngularJS\ped{G}}:} \\
		\begin{itemize}
			\item \textit{Documentazione generica:} \url{https://docs.angularjs.org/guide}
			\item \textit{Documentazione servizio \$http:} \url{https://docs.angularjs.org/api/ng/service/\$http}
			\item \textit{Documentazione servizio \$location:} \url{https://docs.angularjs.org/api/ng/service/\$location}
			\item \textit{Documentazione servizio \$windows:} \url{https://docs.angularjs.org/api/ng/service/\$window}
			\item \textit{Documentazione servizio \$ruoteParams:} \url{https://docs.angularjs.org/api/ngRoute/service/\$routeParams}
			\item \textit{Documentazione servizio \$q:} \url{https://docs.angularjs.org/api/ng/service/\$q}
		\end{itemize}
	\item \textbf{Documentazione del \textit{framework\ped{G}} Material for Angular:} \\
		\url{https://material.angularjs.org/latest/}
	\item \textbf{Documentazione del \textit{framework\ped{G}} Chart.js} \\
		\url{http://www.chartjs.org/docs/}
	\item \textbf{Documentazione del \textit{wrapper\ped{G}} Angles.js} \\
		\url{https://github.com/gonewandering/angles}
	\item \textbf{Documentazione del \textit{framework\ped{G}} TextAngular.js} \\
		\url{https://github.com/fraywing/textAngular/wiki/textAngular-Docs-v1.1.x}
	\item \textbf{Guida all'utilizzo della direttiva \textit{ng-file-upload}} \\
		\url{https://github.com/danialfarid/ng-file-upload}
	\item \textbf{Documentazione di \textit{jison\ped{G}} per la definizione della grammatica di QML} \\
		\url{http://zaa.ch/jison/docs/}
\end{itemize}
\newpage
\section{Requisiti di Sistema}
QuizziPedia è un'applicazione web, pertanto è necessario che il dispositivo di utilizzo abbia una connessione a Internet. Non sono richieste particolari prestazioni per quanto riguarda l'hardware del dispositivo.
\subsection{Dispositivi supportati}
Il software \progetto{} è compatibile con i seguenti sistemi operativi desktop: \textit{Ubuntu\ped{G}} 16 LTS, \textit{OS\ped{G}} X El Capitan 10 e \textit{Windows\ped{G}} 10 e con i seguenti sistemi operativi mobile: \textit{Android\ped{G}} 6.0 Marshmallow, \textit{IOS\ped{G}} 9 e \textit{Windows 10 Mobile\ped{G}}.
\subsection{Browser supportati}
Il software \progetto{} supporta i browser: \textit{Google Chrome\ped{G}} 50, \textit{Mozilla Firefox\ped{G}} 46, \textit{Microsoft Edge\ped{G}} 25 e \textit{Opera\ped{G}} 37.
\newpage
\section{Configurazione e modalità di utilizzo}
\subsection{Configurazione dell'ambiente di lavoro}
Per l'utilizzo del software è necessaria l'installazione della piattaforma \textit{Node.js\ped{G}}, del Node \textit{package\ped{G}} manager \textit{npm\ped{G}} e del sistema software di controllo di versione distribuito \textit{Git\ped{G}}. 
Inoltre nell'ambiente di lavoro in cui verrà eseguito QuizziPedia dovrà essere installato Python 2.7 per il corretto funzionamento dei test, eseguendo il comando da terminale:\\
\\
\centerline{\texttt{npm install -g node-gyp}}\\
\\
e installando, dipendentemente dal sistema operativo in uso, i seguenti pacchetti:
\begin{itemize}
	\item \textbf{UNIX}
	\begin{itemize}
		\item Python (v2.7 raccomandata, v3.x.x non è supportata);
		\item \texttt{make};
		\item Un compilatore C/C++, come GCC.
	\end{itemize}
	\item \textbf{Mac OS X}
	\begin{itemize}
		\item Python (v2.7 raccomandata, v3.x.x non è supportata);
		\item Xcode: è necessario installare il Command Line Tools tramite Xcode. Lo si può trovare nel menu Xcode -> Preferences -> Downloads;
		\item Questo passo installerà gcc e il relativo pacchetto di programmi contiene il \texttt{make}.
	\end{itemize}
	\item \textbf{Windows 10}
	\begin{itemize}
		\item Python (v2.7.10 raccomandato, v3.x.x non è supportata): assicurarsi che la variabile d'ambiente PYTHON abbia il valore: \verb|\path\to\python.exe|;
		\item Installare l'ultima versione di \textit{npm\ped{G}};
		\item Installare Visual Studio Community 2015 Edition;
		\item Impostare la variabile d'ambiente \verb|GYP_MSVS_VERSION=2015|;
		\item Avviare il prompt dei comandi come Amministratore ed eseguire il comando: \texttt{npm install};
		\item Per i sistemi a 64-bit è necessario anche Windows 7 64-bit SDK;
		\item Potrebbe essere necessario eseguire uno dei seguenti comandi se WindowsSDKDir lamenta di non essere impostato:
	\end{itemize}
\end{itemize}
	\begin{center}
		\verb|call "C:\Program Files\Microsoft SDKs\Windows\v7.1\bin\Setenv.cmd"/Release/x86|\\
		\verb|call "C:\Program Files\Microsoft SDKs\Windows\v7.1\bin\Setenv.cmd"/Release/x64|\\
	\end{center}
La guida originale per l'installazione del pacchetto node-gyp è visibile sul link presentato all'interno della sezione Riferimenti -> Informativi.\\
A questo punto è possibile scaricare il software presente nel seguente link:\\ 
\url{https://github.com/TheFellowshipOfTheCode/QuizziPedia}

\subsection{Configurazione del database}
Affinché l'applicazione funzioni correttamente è necessaria la configurazione del database. Per lo sviluppo del progetto è stata utilizzata la piattaforma mLab (\url{https://mlab.com}), che offre la possibilità di creare database Mongo. Dopo aver effettuato l'iscrizione e l'autenticazione al servizio, è necessaria la creazione di un database e di un utente che vi può accedere. Il nome del database deve essere \texttt{quizzipedia} e le collections che devono essere inserite al suo interno sono:
\begin{itemize}
	\item \texttt{questions}
	\item \texttt{quizzes}
	\item \texttt{topics}
	\item \texttt{users}
	\item \texttt{Variables}
\end{itemize}
 La collection \texttt{Variables} contiene le keywords per la traduzione di tutte le parole chiave presenti nel software. Per la traduzione in italiano scaricare il file JSON presente all'indirizzo:\\ 
 \url{https://github.com/TheFellowshipOfTheCode/Documenti/blob/master/Variables_it.json}\\
 mentre per quella in inglese all'indirizzo:\\ 
 \url{https://github.com/TheFellowshipOfTheCode/Documenti/blob/master/Variables_eng.json}.\\ 
 Per ultimare la configurazione del database configurare il file \texttt{loginToMongoLab} presente nel percorso QuizziPedia/Back-End/Config:
 \begin{lstlisting}[language=Java,firstnumber=1]
 var login="nome_utente_database";
 var password="password_utente_database";
 var url="url_database_quizzipedia";
 var database="quizzipedia";
 
 module.exports.login= login;
 module.exports.password= password;
 module.exports.url= url;
 module.exports.database= database;
 
 \end{lstlisting}
 
\subsection{Avvio dell'applicazione}
Da terminale eseguire i seguenti comandi all'interno della cartella QuizziPedia:\\
\\
\centerline{\texttt{npm install}}\\
\\
\centerline{\texttt{npm start}}\\
\\
A questo punto il software è avviabile attraverso un \textit{browser\ped{G}} sulla porta \textit{localhost\ped{G}} indicata dal terminale.

\subsection{Modalità di utilizzo}
Vi sono tre modalità di utilizzo del software, ognuna delle quali ha differenti privilegi:
\begin{itemize}
	\item \textbf{Utente non autenticato}: non è necessaria l'iscrizione all'applicazione ed è possibile solo effettuare la modalità allenamento e ricercare utenti e questionari;
	\item \textbf{Utente autenticato}: è necessaria l'iscrizione attraverso l'apposita funzionalità visibile nell'home page, compilandone tutti i campi dati. Una volta effettuata l'autenticazione, è possibile accedere a tutte le funzionalità che l'applicazione offre tranne la creazione di questionari;
	\item \textbf{Utente autenticato pro}: il software permette di effettuare un upgrade del proprio account all'interno della funzionalità Gestione profilo. Un utente autenticato pro può accedere ad ogni funzionalità dell'applicazione. 
\end{itemize}


\newpage
\section{Librerie e frameworks utilizzati}
\subsection{Front-End}
Nella parte Front-End dell'applicazione sono state utilizzate le seguenti librerie e \textit{frameworks\ped{G}}:
\begin{itemize}
	\item \texttt{MaterialAngular.js}: \textit{framework\ped{G}} per lo sviluppo dei componenti grafici dell'applicazione;
	\item \texttt{Charts.js}: libreria per lo sviluppo dei grafici per la visualizzazione delle statistiche
	utente;
	\item \texttt{Angles.js}: libreria necessaria per integrare la libreria Charts.js all'interno
	dell'ambiente \textit{Angular\ped{G}};
	\item \texttt{Ace}: editor di codice incorporabile scritto in JavaScript.
	\item
	\texttt{angular-number-picker}: direttiva usata per la raccolta dei numeri tramite pulsante su / giù, invece di digitarli;
	\item
	\texttt{AngularCSS}: libreria che ottimizza il livello di presentazione delle single page application iniettando i fogli di stile in modo dinamico in base alle esigenze;
	\item 
	\texttt{AngularUI}: libreria che permette di utilizzare direttive Bootstrap all'interno dell'ambiente Angular;
	\item
	\texttt{Drag and Drop for AngularJS}: libreria che implementa la funzionalità jQuery UI Drag and Drop nell'ambiente Angular;
	\item \texttt{Jison}: generatore di parser JavaScript.
\end{itemize}
\subsection{Back-End}
Nella parte Back-End dell'applicazione sono state utilizzate le seguenti librerie e \textit{frameworks\ped{G}}:
\begin{itemize}
	\item \texttt{Express}: \textit{framework\ped{G}} Web di routing e \textit{middleware\ped{G}}, con funzionalità sua propria minima: un'applicazione Express è essenzialmente una serie di chiamate a funzioni \textit{middleware\ped{G}};
	\item \texttt{Passport}: \textit{middleware\ped{G}} per l'autenticazione in ambiente \textit{Node.js\ped{G}}.
\end{itemize}

%Generazione delle variabili che andranno a sostituire quelle del template `HomePage.tex`
\newcommand{\documento}{\G}
\newcommand{\nomedocumentofisico}{Glossario.pdf}
\newcommand{\redazione}{\SM}
\newcommand{\verifica}{\FB}
\newcommand{\approvazione}{Da Approvare}
\newcommand{\uso}{Esterno}
\newcommand{\destinateTo}{\TV \\ & \RC \\ & \proponente \\ & \gruppo}
\newcommand{\datacreazione}{24 Dicembre 2015}
\newcommand{\datamodifica}{24 Dicembre 2015}
\newcommand{\stato}{Sviluppo}

\def\TABELLE{false}	%abilita - disabilita l'indice delle tabelle
\def\FIGURE{false} 	%abilita - disabilita l'indice delle figure

\documentclass[a4paper,11pt]{article}

%***IMPORTAZIONE PACKAGE***
\usepackage{ifthen}
\usepackage[italian]{babel}
\usepackage[utf8]{inputenc}
\usepackage[T1]{fontenc}
\usepackage{float}
\usepackage{chapterbib}
\usepackage{graphicx}
\usepackage[a4paper,top=2.5cm,bottom=2.5cm,left=2.5cm,right=2.5cm]{geometry}
\usepackage[colorlinks=true, urlcolor=black, citecolor=black, linkcolor=black]{hyperref}
\usepackage{booktabs}
\usepackage{fancyhdr}
\usepackage{totpages}
\usepackage{tabularx, array}
\usepackage{dcolumn}
\usepackage{epstopdf}
\usepackage{booktabs}
\usepackage{fancyhdr}
\usepackage{longtable}
\usepackage{calc}
\usepackage{datatool}
\usepackage[bottom]{footmisc}
\usepackage{listings} 
\usepackage{textcomp}
\usepackage{titlesec}
\usepackage{rotating} 
\usepackage{multirow}
\usepackage{placeins}
\usepackage{color}
\usepackage[table,usenames,dvipsnames]{xcolor}
\usepackage{hyperref}
\usepackage{makecell}
\usepackage{hyperref}


%***STILE PAGINA***
\pagestyle{fancy}
%no indentazione paragrafo
\setlength{\parindent}{0pt}

%***INTESTAZIONE***
\lhead{\Large{\progetto} \\ \footnotesize{\documento}}
\rhead{\includegraphics[keepaspectratio = true, width = 25px] {../../Template/icone/LogoGruppo.png}}
\renewcommand{\headrulewidth}{0.4pt}  %Linea sotto l'intestazione

%***PIÈ DI PAGINA***
\lfoot{\textit{\gruppoLink}\\ \footnotesize{\email}}
\rfoot{\thepage} %per le prime pagine: mostra solo il numero romano
\cfoot{}
\renewcommand{\footrulewidth}{0.4pt}   %Linea sopra il piè di pagina

%***INSERIMENTO DI NUOVE SOTTOSEZIONI
\setcounter{secnumdepth}{7}		% mostra nel documento fino al livello 8 (1.2.3.4.5.6.7.8)
\setcounter{tocdepth}{7}			% mostra nell'indice fino al livello 8 (1.2.3.4.5.6.7.8)

%***LA SOTTOSEZIONE PARAGRAPH VIENE VISUALIZZATA COME UNA SECTION
\titleformat{\paragraph}{\normalfont\normalsize\bfseries}{\theparagraph}{1em}{}
\titlespacing*{\paragraph}{0pt}{3.25ex plus 1ex minus .2ex}{1.5ex plus .2ex}

\titleformat{\subparagraph}{\normalfont\normalsize\bfseries}{\thesubparagraph}{1em}{}
\titlespacing*{\subparagraph}{0pt}{3.25ex plus 1ex minus .2ex}{1.5ex plus .2ex}

\makeatletter
\newcounter{subsubparagraph}[subparagraph]
\renewcommand\thesubsubparagraph{%
  \thesubparagraph.\@arabic\c@subsubparagraph}
\newcommand\subsubparagraph{%
  \@startsection{subsubparagraph}    % counter
    {6}                              % level
    {\parindent}                     % indent
    {3.25ex \@plus 1ex \@minus .2ex} % beforeskip
    {0.75em}                           % afterskip
    {\normalfont\normalsize\bfseries}}
\newcommand\l@subsubparagraph{\@dottedtocline{6}{10em}{5.5em}} %gestione dell'indice
\newcommand{\subsubparagraphmark}[1]{}
\makeatother

\makeatletter
\newcounter{subsubsubparagraph}[subsubparagraph]
\renewcommand\thesubsubsubparagraph{%
  \thesubsubparagraph.\@arabic\c@subsubsubparagraph}
\newcommand\subsubsubparagraph{%
  \@startsection{subsubsubparagraph}    % counter
    {7}                              % level
    {\parindent}                     % indent
    {3.25ex \@plus 1ex \@minus .2ex} % beforeskip
    {0.75em}                           % afterskip
    {\normalfont\normalsize\bfseries}}
\newcommand\l@subsubsubparagraph{\@dottedtocline{7}{10em}{6.5em}} %gestione dell'indice
\newcommand{\subsubsubparagraphmark}[1]{}
\makeatother
%Generali
\newcommand{\progetto}{QuizziPedia}
\newcommand{\gruppo}{TheFellowshipOfTheCode}
\newcommand{\gruppoLink}{\href{http://thefellowshipofthecode.github.io/}{TheFellowshipOfTheCode}}
\newcommand{\email}{\href{mailto:thefellowshipofthecode@gmail.com}{thefellowshipofthecode@gmail.com}}

%Documenti
\newcommand{\AdR}{Analisi dei Requisiti}
\newcommand{\NdP}{Norme di Progetto}
\newcommand{\PdP}{Piano di Progetto}
\newcommand{\SdF}{Studio di Fattibilità}
\newcommand{\PdQ}{Piano di Qualifica}
\newcommand{\VI}{Verbale Interno}
\newcommand{\VE}{Verbale Esterno}
\newcommand{\ST}{Specifica Tecnica}
\newcommand{\DDP}{Definizione di Prodotto}
\newcommand{\MU}{Manuale Utente}
\newcommand{\G}{Glossario}
\newcommand{\LdP}{Lettera di Presentazione}
\newcommand{\NdPv}{NormeDiProgetto\_v\_1\_0\_0}
\newcommand{\PdPv}{PianoDiProgetto\_v\_1\_0\_0}
\newcommand{\PdQv}{PianoDiQualifica\_v\_1\_0\_0}
\newcommand{\SdFv}{StudioDiFattibilità\_v\_1\_0\_0}

%Componenti del gruppo
\newcommand{\AF}{Alberto Ferrara}
\newcommand{\SM}{Simone Magagna}
\newcommand{\FB}{Franco Berton}
\newcommand{\MP}{Marco Prelaz}
\newcommand{\MV}{Mattia Varotto}
\newcommand{\GN}{Matteo Gnoato}
\newcommand{\GR}{Matteo Granzotto}

%Ruoli
\newcommand{\RdP}{Responsabile di Progetto}
\newcommand{\Res}{Responsabile}
\newcommand{\Amm}{Amministratore}
\newcommand{\Ver}{Verificatore}
\newcommand{\Prog}{Progettista}
\newcommand{\Progr}{Programmatore}
\newcommand{\Ana}{Analista}
\newcommand{\RdPs}{Responsabili di Progetto}
\newcommand{\Ress}{Responsabile}
\newcommand{\Amms}{Amministratori}
\newcommand{\Vers}{Verificatori}
\newcommand{\Progs}{Progettisti}
\newcommand{\Progrs}{Programmatori}
\newcommand{\Anas}{Analisti}

%Professori e proponente
\newcommand{\TV}{Prof. Tullio Vardanega}
\newcommand{\RC}{Prof. Riccardo Cardin}
\newcommand{\ZU}{Zucchetti S.P.A.}
\newcommand{\proponente}{Zucchetti S.P.A.}

\newcommand{\diaryEntry}[5]{#2 & \emph{#4} & #3 & #5 & #1\\ \hline}

%comando per una nuova riga nella tabella del diario delle modifiche
\newcommand{\specialcell}[2][c]{%
	\begin{tabular}[#1]{@{}c@{}}#2\end{tabular}}

\renewcommand*\sectionmark[1]{\markboth{#1}{}}
\renewcommand*\subsectionmark[1]{\markright{#1}}

%Pediodi di lavoro 
\newcommand{\AR}{Analisi dei Requisiti}
\newcommand{\AD}{Analisi dei Requisiti in Dettaglio}
\newcommand{\PA}{Progettazione Architetturale}
\newcommand{\PD}{Progettazione di Dettaglio}
\newcommand{\CO}{Codifica}
\newcommand{\VV}{Verifica e Validazione}

% Revisioni
\newcommand{\RR}{Revisione dei Requisiti}
\newcommand{\RP}{Revisione di Progettazione}
\newcommand{\RQ}{Revisione di Qualifica}
\newcommand{\RA}{Revisione di Accettazione}

% Comandi analisi dei requisiti
\newcommand{\uau}{utente autenticato}
\newcommand{\uaus}{utenti autenticati}
\newcommand{\uaupro}{utente autenticato pro}
\newcommand{\uauspro}{utenti autenticati pro}

\newcommand{\myincludegraphics}[2][]{%
	\setbox0=\hbox{\phantom{X}}%
	\vtop{
		\hbox{\phantom{X}}
		\vskip-\ht0
		\hbox{\includegraphics[#1]{#2}}}}


\newcommand{\modificheuno} 
{	
	0.0.3 & Inseriti i primi riferimenti informativi & \specialcell[t]{\GN \\ \prog} & 2016-02-26
	\\\midrule
	0.0.2 & Stesura scopo del documento, scopo del prodotto , riferimenti normativi & \specialcell[t]{\GN \\ \prog} & 2016-02-26
	\\\midrule
	0.0.1 & Creato template documento & \specialcell[t]{\GN \\ \prog} & 2016-02-26
	\\\midrule
}
\newcommand{\modifichedue}
{
}
%Serve per vedere i paragraph nell'indice
\setcounter{secnumdepth}{7}
\setcounter{tocdepth}{7}
\begin{document}

%inclusione template HomePage
\begin{center}

%\includegraphics[width=1em]{../../../Template/icone/LogoGruppo.png}
\begin{large} \textbf{\gruppoLink} \end{large}
%\includegraphics[width=1em]{../../../Template/icone/LogoGruppo.png}
\vspace{0.2em}

\hrule
\vspace{3em}

\includegraphics[keepaspectratio = true, width=8cm]{../../../Template/icone/LogoGruppo.png}

%Prima pagina senza intestazione né piè di pagina	
\thispagestyle{empty}

%Le informazioni del documento sono ancorate a fine pagina
\vfill

%Copertina
\begin{center} 
  \begin{Huge}
  {\fontsize{15mm}{20mm}\selectfont \progetto} 
  \end{Huge}
\end{center}

\begin{Huge} \documento \end{Huge}

\begin{center}
\textbf{Informazioni sul documento} \\ \vspace{2em}
\small
\begin{tabular}{r|l}
	\textbf{Nome Documento} & \nomedocumentofisico \\
	\textbf{Versione}	& 1\\
	\textbf{Data di Creazione} & \datacreazione\\
	\textbf{Data ultima modifica} & \datamodifica\\
	\textbf{Stato} & \stato \\
	\textbf{Redazione}	& \redazione\\
	\textbf{Verifica}	& \verifica\\
	\textbf{Approvazione}	& \approvazione\\
	\textbf{Uso}  & \uso\\
	\textbf{Distribuzione} & \gruppo \\
	\textbf{Destinato a}  &  \destinateTo \\
	\textbf{Email di riferimento} & \email
\end{tabular}
\end{center}

\normalsize
%Sommario
\textbf{Sommario\\} 
Documento contenente le norme di progetto che il gruppo \textit{\gruppo} seguirà durante tutte le fasi di realizzazione del prodotto \textit{\progetto}.

%\vfill %cosa fa?
\end{center}
\clearpage


%Registro delle modifiche e indice 
%si usa la numerazione romana per gli indici e la tabella delle modifiche
\pagenumbering{Roman}
\newpage
%***REGISTRO DELLE MODIFICHE***
%Vari comandi per la struttura della tabella, NON MODIFICARE!
\begin{center}
	\Large{\textbf{Registro delle modifiche}}
	\\\vspace{0.5cm}
	\normalsize
	\begin{tabularx}{\textwidth}{cXcc}
		\textbf{Versione} & \textbf{Descrizione} & \textbf{Autore e Ruolo} & \textbf{Data} \\\toprule
		\modifiche
		\bottomrule
\end{tabularx}
\end{center}
\newpage
%Inserisce il link all'indice
%\addcontentsline{toc}{section}{Indice}
\tableofcontents
\clearpage 

%Se è stata impostata a true la variabile per la lista delle tabelle, la mostra
\ifthenelse{\equal{\TABELLE}{true}} 
{\listoftables \newpage}{}

%Se è stata impostata a true la variabile per la lista delle figure, la mostra
\ifthenelse{\equal{\FIGURE}{true}}
{\listoffigures \newpage}{}

%Da qui comincia la numerazione normale
\pagenumbering{arabic}

%Imposta il formato di visualizzazione
\rfoot{\thepage~di~\pageref{TotPages}}
\end{document}
\appendix
%\newpage
\section{Organigramma}
\subsection{Redazione}
\begin{table}[htbp]
	\begin{center}
		\setlength{\extrarowheight}{\jot}
		\begin{tabular}{|c|c|p{6cm}|}
			\hline
			\textbf{Nominativo} & \textbf{Data di redazione} & \textbf{Firma} \\[1ex]
			\hline
			\GR & 2016-01-07 & \myincludegraphics{immagini/Firme/MGR.png} \\[1ex]
			\hline
		\end{tabular}
	\end{center}
	\caption{Redazione}
\end{table}

\subsection{Approvazione}
\begin{table}[htbp]
	\begin{center}
		\setlength{\extrarowheight}{\jot}
		\begin{tabular}{|c|c|p{5cm}|}
			\hline
			\textbf{Nominativo}     & \textbf{Data di approvazione} & \textbf{Firma}  \\[1ex]
			\hline
			\GR		& 2016-01-07					& \myincludegraphics{immagini/Firme/MGR.png}			\\[1ex]
			\hline
			\TV	&								&			\\[1ex]
			\hline
		\end{tabular}
	\end{center}
	\caption{Approvazione}
\end{table}

\subsection{Accettazione dei componenti}
\begin{table}[htbp]
	\begin{center}
		\setlength{\extrarowheight}{\jot}
		\begin{tabular}{|c|c|p{6cm}|}
			\hline
			\textbf{Nominativo} & \textbf{Data di accettazione} & \textbf{Firma} \\[1ex]
			\hline
			\GR	&	2016-01-07	& \myincludegraphics{immagini/Firme/MGR.png}	\\[1ex]
			\hline
			\GN		&	2016-01-07	& \myincludegraphics{immagini/Firme/MGN.png}		\\[1ex]
			\hline
			\AF		&	2016-01-07	& \myincludegraphics{immagini/Firme/AF.png}		\\[1ex]
			\hline
			\FB		&	2016-01-07	& \myincludegraphics{immagini/Firme/FB.png}		\\[1ex]
			\hline
			\MP		&	2016-01-07	& \myincludegraphics{immagini/Firme/MP.png}		\\[1ex]
			\hline
			\SM		&	2016-01-07	& \myincludegraphics{immagini/Firme/SM.png}		\\[1ex]
			\hline
			\MV		&	2016-01-07	& \myincludegraphics{immagini/Firme/MV.png}		\\[1ex]
			\hline
		\end{tabular}
	\end{center}
	\caption{Accettazione}
\end{table}

\subsection{Componenti}
\begin{table}[H]
	\begin{center}
		\setlength{\extrarowheight}{\jot}
		\begin{tabular}{|c|c|p{5cm}|p{4.3cm}|}
			\hline
			\textbf{Nominativo} & \textbf{Matricola} & \raggedright \textbf{Indirizzo di posta elettronica} & \textbf{Ruoli} \\[1ex]
			\hline
	 		\GR	& 1051540	& \href{mailto:granzotto.matteo@gmail.com}{granzotto.matteo@gmail.com} 	& \Res, \Ana 	\\[1ex]
			\hline
			\GN		& 1051873	& \href{mailto:gnoatomatteo@gmail.com}{gnoatomatteo@gmail.com} 			& \Res, \Ana 	\\[1ex]
			\hline
			\AF		& 1049378	& \href{mailto:albertoferrara92@gmail.com}{albertoferrara92@gmail.com} 	& \Amm, \Ana 	\\[1ex]
			\hline
			\FB 		& 1052574	& \href{mailto:franco.berton93@gmail.com}{franco.berton93@gmail.com} 	& \Amm, \Ana	\\[1ex]
			\hline
			\MP		& 1047343	& \href{mailto:marcomak91@hotmail.it}{marcomak91@hotmail.it} 			& \Ver, \Ana	\\[1ex]
			\hline
			\SM		& 1009467	& \href{mailto:simone.magagna91@gmail.com}{simone.magagna91@gmail.com} 	& \Ver, \Ana, \Amm 	\\[1ex]
			\hline
			\MV		& 1051619	& \href{mailto:varots93@hotmail.it}{varots93@hotmail.it} & \Ver, \Ana, \Amm \\[1ex]
			\hline	
		\end{tabular}
	\end{center}
	\caption{Componenti}
\end{table}



\end{document}